\emph{Setting and two preliminary facts.} Everything is in the module and with the notation of Theorem \claimref{synthesis.positive-structure-is-twistor-form-at-imaginary-time}; the check realises it as three fermionic modes in the Jordan-Wigner basis over the Gaussian rationals. Write $A^{\sharp}$ for the $\Sigma$-adjoint. By (iii) of that theorem every real vector is $\Sigma$-anti-Hermitian, so for the real bivector $e_{1}e_{2}$, $(e_{1}e_{2})^{\sharp} = e_{2}^{\sharp}e_{1}^{\sharp} = e_{2}e_{1} = -e_{1}e_{2}$, and $K^{\sharp} = -K$. Write $J_{m}$ for the involution equal to $+1$ on $L(m)$ and $-1$ on $L(\bar m)$, so that $h_{m} = \Sigma(\cdot, J_{m}\cdot)$; the $h_{m}$-adjoint is $A^{\dagger_{m}} = J_{m}A^{\sharp}J_{m}$.

\emph{Covariance.} For a rotor $S$ of $\Spin(2,4)$ acting on points by $g$ (so that $SX(m)S^{-1}$ is a nonzero multiple of $X(gm)$, Proposition \claimref{one-relation.spacetime-manifold-vacuum-lines} extended complex-linearly), $S L(m) = L(gm)$, hence $J_{gm} = SJ_{m}S^{-1}$ and, $\Sigma$ being invariant, $h_{gm}(\psi, \varphi) = h_{m}(S^{-1}\psi, S^{-1}\varphi)$. The boosts $S_{\lambda}$ and the dilations $\exp(\tfrac{\mu}{2}e_{+}e_{-})$ commute with $N$: they fix $\nO$, $\ninf$ up to reciprocal factors, and $N = \tfrac12\ninf\nO$ (the check confirms $S_{\lambda}$ fixes $a$, $a^{*}$ and that $e_{+}e_{-}$ commutes with $N$). So for $g$ generated by them, $\langle N\rangle$ in the state $\psi$ read with $h_{gm}$ equals $\langle N\rangle$ in $S^{-1}\psi$ read with $h_{m}$.

\emph{(i)} $W_{\theta}$ is $h_{m}$-unitary for all $\theta$ if and only if $iK$ is $h_{m}$-anti-Hermitian. As $\sharp$ is antilinear, $(iK)^{\sharp} = -iK^{\sharp} = iK$, so $(iK)^{\dagger_{m}} = iJ_{m}KJ_{m}$, and the condition is $J_{m}KJ_{m} = -K$: $K$ anticommutes with $J_{m}$, equivalently $K$ maps $L(m)$ into $L(\bar m)$ and $L(\bar m)$ into $L(m)$.
To read this geometrically, let $P_{m}$ be the real plane spanned by $\operatorname{Re}X(m)$ and $\operatorname{Im}X(m)$. At $m = ie_{1}$, $X = e_{+} + ie_{1}$ and $J_{ie_{1}} = ie_{1}e_{+}$ is $i$ times the unit bivector of $P_{ie_{1}} = \operatorname{span}(e_{1}, e_{+})$, a plane on which the metric is positive definite. The forward tube is one orbit of the identity component of $\Spin(2,4)$, and by covariance $J_{m} = i I_{m}$ with $I_{m}$ a unit bivector of $P_{m}$ (square $-1$) for every tube point. Two 2-blades $A$, $B$ anticommute exactly when the symmetric part $\tfrac12(AB + BA) = A \cdot B + A \wedge B$ of their product vanishes, that is when their planes share a line ($A \wedge B = 0$) and $A \cdot B = 0$. Apply this to $A = e_{1}e_{2}$ and $B = I_{m}$. Write $m = x + iy$; then $\operatorname{Re}X(m) = \nO + x - \tfrac12\operatorname{Re}Q(m)\ninf$ and $\operatorname{Im}X(m) = y - \tfrac12\operatorname{Im}Q(m)\ninf$. A vector $\alpha\operatorname{Re}X + \beta\operatorname{Im}X$ in $\operatorname{span}(e_{1}, e_{2})$ has no $\nO$-component, so $\alpha = 0$, and then $\operatorname{Im}Q(m) = 2x \cdot y = 0$ and $y \in \operatorname{span}(e_{1}, e_{2})$; $y$ is future timelike because $m$ is in the tube. With $c = \operatorname{Re}X$ and $d = y$, $(e_{1} \wedge e_{2}) \cdot (c \wedge d)$ is a nonzero multiple of $(e_{1} \cdot d)(e_{2} \cdot c) - (e_{1} \cdot c)(e_{2} \cdot d) = (e_{1} \cdot y)(e_{2} \cdot x) - (e_{1} \cdot x)(e_{2} \cdot y)$, which vanishes exactly when the component $x_{\parallel}$ of $x$ in the boost plane is parallel to $y$. Together with $x_{\parallel} \cdot y = x \cdot y = 0$ and $y$ timelike this forces $x_{\parallel} = 0$. Conversely, for $x \in \operatorname{span}(e_{3}, e_{4})$ and $y$ in the boost plane both conditions hold. So the condition is $m = x_{\perp} + i\rho u$ as stated. The wedge reflection sends $x + iy$ to $R(x) - iR(y)$, which equals $x + iy$ exactly when $R(x) = x$ and $R(y) = -y$, that is $x \in \operatorname{span}(e_{3}, e_{4})$ and $y \in \operatorname{span}(e_{1}, e_{2})$: the same set. It has real dimension two ($x_{\perp}$) plus two ($\rho$, $\eta$). The check builds $h_{m}$ from the two lines at four compatible points (including $ie_{1}$, a boosted point, and points with transverse real part), confirms it is positive definite, that $iK$ is $h_{m}$-anti-Hermitian on the half and that $e_{1}e_{2}$ swaps $L(m)$ and $L(\bar m)$, and at five tube points off the set (real part with a boost-plane component, imaginary part leaving the boost plane) confirms that $h_{m}$ is positive and $iK$ is not anti-Hermitian; it also confirms $h_{ie_{1}} = h$ on the half and the reflection test at all nine points.

\emph{(ii)} $(e_{1}e_{2})^{2} = -e_{1}^{2}e_{2}^{2} = 1$, so $W_{\theta} = \cos\tfrac{\theta}{2} + i\sin\tfrac{\theta}{2}\,e_{1}e_{2}$, and $e_{1}e_{2}$, being even and not $\pm 1$ on a half, has the eigenvalues $\pm 1$ there; the check finds each eigenspace two-dimensional. Hence the eigenvalues $e^{\pm i\theta/2}$, $W_{2\pi} = \cos\pi = -1$ and $W_{4\pi} = 1$. The check confirms the rotor at the rational half angle $(\cos, \sin) = (\tfrac35, \tfrac45)$ is $h$-unitary and composes by the double-angle law, and with $1 + ie_{1}e_{2} = \sqrt{2}\,W_{\pi/2}$ that $(1 + ie_{1}e_{2})^{2} = 2ie_{1}e_{2}$ and $(1 + ie_{1}e_{2})^{4} = -4$, the exact form of $W_{\pi} = ie_{1}e_{2}$ and $W_{2\pi} = -1$.

\emph{(iii)} Let $B = e_{1}e_{2}$ and $U = \cos\phi + i\sin\phi\,B$ with $\phi = \pi/4$. $B$ anticommutes with $e_{1}$ and $e_{2}$ and commutes with $e_{3}, e_{4}, e_{\pm}$, so $U$ fixes $\nO$, $\ninf$, $e_{3}$, $e_{4}$, and $Ue_{1}U^{-1} = U^{2}e_{1} = iBe_{1} = -ie_{2}$, $Ue_{2}U^{-1} = iBe_{2} = -ie_{1}$. Hence $U(i\rho u)U^{-1} = \rho(\cosh\eta\,e_{2} + \sinh\eta\,e_{1})$, and since $Q$ is invariant, $UX(m)U^{-1} = X(x(m))$. The point $x(m) - x_{\perp} = \rho(\sinh\eta\,e_{1} + \cosh\eta\,e_{2})$ has $Q = -\rho^{2} < 0$ and positive $e_{2}$-component; conversely every right-wedge point has a unique such form with $\rho > 0$ and real $\eta$, so the map is a bijection. $U$ commutes with $S_{\lambda}$ (both functions of $B$), with $e_{+}e_{-}$ and with the translations along the bifurcation plane (whose rotors are built from $\ninf$ and $e_{3}$, $e_{4}$, all commuting with $B$), which gives the equivariance. The tangent of the boost orbit $\lambda \mapsto \rho(\sinh(\eta + \lambda)e_{1} + \cosh(\eta + \lambda)e_{2})$ at $\lambda = 0$ is $\rho u$, and $\rho$ is the distance from the bifurcation plane. At $\rho = 1$, $\eta = 0$, $x_{\perp} = 0$: $m = ie_{1}$, $x(m) = e_{2}$. The boost changes $\eta$ and nothing else, so the compatible family modulo boosts is the set of pairs $(x_{\perp}, \rho)$ with $\rho > 0$, an upper half space; the dilation by $\lambda$ sends $(x_{\perp}, \rho)$ to $(\lambda x_{\perp}, \lambda\rho)$ and the translation by $c \in \operatorname{span}(e_{3}, e_{4})$ sends it to $(x_{\perp} + c, \rho)$, both commuting with $K$, so together they act transitively on it. The check verifies $UX(m)U^{-1} = X(x(m))$ exactly at three compatible points with $(\rho, \cosh\eta, \sinh\eta) = (1, 1, 0), (2, \tfrac54, \tfrac34), (\tfrac13, \tfrac{17}{8}, -\tfrac{15}{8})$ and a transverse part $\tfrac27 e_{3}$, and that $S_{\lambda}$ fixes $\nO$ and commutes with $N$.

\emph{(iv)} Let $c \in \operatorname{span}(e_{3}, e_{4})$. Since $c$ is orthogonal to $\nO$ and $\ninf$ and $\ninf$ is null, $\{a_{c}, a_{c}^{*}\} = \tfrac12\{\nO, \ninf\} = 1$. $S_{\lambda}$ commutes with $\nO$, $\ninf$ and $c$ (the bivector $e_{1}e_{2}$ commutes with each), so $a_{c}$, $a_{c}^{*}$ are fixed by the automorphisms $\operatorname{Ad}S_{\lambda}$ and by their entire extension in $\lambda$. For a functional $\omega$ satisfying the KMS condition $\omega(A\,\alpha_{i\beta}(B)) = \omega(BA)$ for the boost automorphisms $\alpha$, take $A = a_{c}^{*}$, $B = a_{c}$: $\omega(a_{c}^{*}a_{c}) = \omega(a_{c}a_{c}^{*})$, and with the anticommutator this is $\langle N_{c}\rangle = \tfrac12$, whatever $\beta$. No positivity is used; by (iv) of the cited theorem the boost is unitary for no positive structure, so the condition is algebraic here.
\emph{Hermiticity fixes the foot.} For $h = h_{ie_{1}}$ the adjoints are $\nO^{\dagger} = \tfrac12\ninf$, $\ninf^{\dagger} = 2\nO$ and $c^{\dagger} = -c$ ($e_{3}, e_{4}$ are unitary of square $-1$), and $Q(c) = -|c|^{2}$. Then $N_{c} = \tfrac12\ninf\nO + \tfrac12\ninf c$ and $N_{c}^{\dagger} = a_{c}^{\dagger}\nO = (\tfrac12\ninf - c + |c|^{2}\nO)\nO = \tfrac12\ninf\nO - c\nO$. As $c$ anticommutes with $\ninf$, $N_{c} = N_{c}^{\dagger}$ if and only if $-\tfrac12 c\ninf = -c\nO$, that is $c(\tfrac12\ninf - \nO) = -ce_{-} = 0$, that is $c = 0$. By covariance (the translation rotor $T_{f}$ along the bifurcation plane sends $X(c)$ to a multiple of $X(c + f)$, fixes $\ninf$, commutes with $K$, and carries $h_{m}$ to $h_{m + f}$; boosts and dilations send $N_{d}$ to $N_{d'}$ with $d' = 0$ exactly when $d = 0$), $N_{c}$ is $h_{m}$-Hermitian for compatible $m$ with foot $x_{\perp}$ exactly when $c = x_{\perp}$.
Now translate by $-c$, which reduces to $c = 0$: let $m$ be compatible with $x_{\perp} = 0$, $m = i\rho' u'$ with $u' = \cosh\eta'\,e_{1} + \sinh\eta'\,e_{2}$, and $x = \rho(\sinh\eta\,e_{1} + \cosh\eta\,e_{2})$. $N$ is $h$-Hermitian, so by covariance it is $h_{m}$-Hermitian. The rotor $g$ (boost by $\eta'$ after dilation by $\rho'$) sends $ie_{1}$ to $m$, so by covariance the occupations on $L(x)$ read with $h_{m}$ are those on $L(x')$ read with $h$, where $x' = (\rho/\rho')(\sinh\delta\,e_{1} + \cosh\delta\,e_{2})$, $\delta = \eta - \eta'$ (Theorem \claimref{synthesis.boost-stationary-incidence-vacua-read-interval}(i): boosts carry lines to lines and fix $N$). In the null coordinates of Theorem \claimref{synthesis.incidence-vacua-thermofield-real-capacity}, $u = t - r = -(\rho/\rho')e^{-\delta}$ and $v = t + r = (\rho/\rho')e^{\delta}$, and by its clause (iii) the occupations of the vacua on the line fill the closed segment between $u^{2}/(1 + u^{2})$ and $v^{2}/(1 + v^{2})$. Every vacuum has occupation $\tfrac12$ exactly when $u^{2} = v^{2} = 1$, that is $\delta = 0$ and $\rho = \rho'$: $m = m(x)$. At $u^{2} = v^{2} = 1$, clause (ii) of that theorem makes $\Omega_{x}$ and $\Psi_{x}$ thermofield doubles at $e^{-\beta} = 1$, $\beta = 0$: the ascent marginal is the trace and its purifying partner, the boost qubit, is maximally entangled with it. The check computes, at three observers and their three tube points, the $h_{m}$ Gram matrices of $1$ and $N$ on $L(x)$ and finds $\operatorname{Gram}(N) = \tfrac12\operatorname{Gram}(1)$ exactly on the diagonal of the pairing $m = m(x)$ and not off it, with $h_{m}$ positive on every line and $N$ $h_{m}$-Hermitian; at $m = 2ie_{1}$ it finds the identity at $x = 2e_{2}$ and its failure at $e_{2}$, $3e_{2}$ and at a boosted point of the orbit of $e_{2}$. At the feet $c = \tfrac12 e_{3}$, $e_{3}$, $2e_{3} - e_{4}$ it confirms that $a_{c}$, $a_{c}^{*}$ are boost invariant with $\{a_{c}, a_{c}^{*}\} = 1$, that $N_{c}$ is $h_{m}$-Hermitian and tracial on $L(x(m))$ for $m = c + i\rho u$ at two scales, that $N = N_{0}$ is not $h_{m}$-Hermitian there, and that $N_{c}$ is not $h_{ie_{1}}$-Hermitian.

\emph{Reading (argued, not part of the statement).} The selected state of the boundary modes is the I-state, KMS for the boost dynamics at the Hawking temperature (Theorem \claimref{hawking.i-state}), which in rapidity units is the period $2\pi$. Two features of such a state have finite-dimensional images here. Its Euclidean continuation is a rotation in imaginary rapidity, and in the Osterwalder-Schrader picture (cited from memory) that rotation is a symmetry of the Euclidean structure, which is condition (i) on the positive structure; and in the Bisognano-Wichmann theorem (cited from memory) the wedge vacuum's modular conjugation is the wedge reflection, whose fixed set is the same family. So a boost-KMS state, to carry one of these positive structures compatibly, reads a tube point $m = x_{\perp} + i\rho u$, which by (iii) is the tube point $m(x)$ of one observer $x$ of the wedge (the foot is the base point of the ascent pair by (iv); the scale $\rho$ the boost data leave open, the dilations commuting with $K$), a quarter turn of the Euclidean circle of radius $\rho$, the midpoint of the half strip $0 < \theta < \pi$ on which Theorem \claimref{synthesis.twistor-form-states-only-off-the-real-slice}(iii) keeps the lines positive; the boost moves the observer along its orbit and the tube point with it, consistent with the absence of a boost-invariant positive structure. By (iv), at that observer's point the incidence line is compatible with the KMS-forced marginal for every vacuum on it, and for no other tube point of the family. For the open problem's capacity question this gives both readings at once and fixes their places, the same way at every compatible tube point (the capacity a state reads does not depend on its tube point, $\Spin(2,4)$ being transitive on the tube and preserving $\Sigma$ and the reality conjugation): the state reads the complex capacity through the definite line $L(m(x))$, available only off the real slice, and the real capacity at the real point $x$, where the boost qubit and the ascent mode form one maximally entangled pair, one $\ln 2$. The ascent mode has boost weight zero, and a KMS state gives a zero-frequency mode the trace; this is the same mechanism by which KMS forces the zero mode of the s-wave array to the trace (Theorem \claimref{synthesis.s-wave-array-is-a-boundary-mode-set}(iii)), whose doublet carries the per-array $\ln 2$ (synthesis.array-ln2-is-zero-mode-doublet-one-tower); the identification of the ascent mode with that zero mode is argued only. The weights $\pm\tfrac12$ of (ii) and the sign at $2\pi$ are, in rapidity units, the half-integer offset and the antiperiodicity of fermionic boundary modes of boost weight one half (premise (Q) reads its offset from that weight); argued. What the module does not fix is the scale $\rho$: dilations act on the family. If the boost flow is normalised by the surface gravity $\kappa$, the Killing field has unit norm at distance $1/\kappa$ from the bifurcation surface, which would make $\rho = 1/\kappa$ (argued; the series' Rindler reading of the horizon is not re-derived); a different candidate is the radius of the bifurcation sphere of the gravitational field ($2M$ for the eternal hole, $1/H$ for de Sitter), the transverse half space being the space of round metrics on that sphere. Either is a scale premise, not derived here. Filed alongside, synthesis.positive-structure-and-natural-scale-actions-do-not-select-the-almost-einstein-scale proves the converse direction: with the boost, the positive structure at $i\rho e_{1}$ fixes the ray of the de Sitter tractor at $\Lambda = 12/\rho^{2}$. The frozen point of positive $\Lambda$ lies in the ascent plane and commutes with $K$, and its stabiliser in the centraliser of $K$ fixes exactly the point of the half space with foot zero and $\rho^{2} = 12/\Lambda$, that is $\rho = 2/H$ in chart units (argued here, from that result). The argued normalisation $\rho = 1/\kappa = 1/H$ agrees with it only when read with the local scale factor $\sigma = 1 - \Lambda Q/12 = 2$ at that distance; in chart length the two differ by a factor $2$, and which normalisation is physical is not decided here. Nothing here assumes a common frame boost for the planes of the tower; the scale of this plane stays a free datum of the state.
